\documentclass{article}
\usepackage[main]{neurips_2026}
\usepackage[utf8]{inputenc}
\usepackage[T1]{fontenc}
\usepackage{hyperref}
\usepackage{url}
\usepackage{booktabs}
\usepackage{amsmath,amssymb}
\usepackage{microtype}
\usepackage{courier}
\usepackage{graphicx}
\usepackage{xcolor}
\input{presentation/captions.tex}
\title{Task Investment to First Place:\\Aggregation Rules and Score Ceilings}
\author{Anonymous Authors}
\begin{document}
\maketitle
\begin{abstract}
Multi-task leaderboards can reward improvements concentrated on selected tasks. We define conditional investment as the minimum number of target task scores that must improve to reach tied first under a specified score envelope, with rivals fixed and no transfer across tasks. Exact optimization on HELM MMLU (79 entries, 57 subjects) and the Open LLM Leaderboard (4,575 entries, six aggregates) compares mean, median, geometric mean and Borda. With selected scores raised to perfection, median costs are respectively 16, 21, 16 and 46 subjects on HELM, and 2, 3, 3 and 6 aggregates on Open LLM. Yet replacing perfection with per-task observed frontiers reverses the mean--median ordering by average minimum cost on both boards; partial uplift can prevent victory altogether. Random-task baselines and task, representation and competitor-pool checks delimit these findings. Aggregation affects the required breadth of score improvement, but its comparison depends on attainable-score assumptions. These are reproducible score-table counterfactuals, not measurements of training expense or contamination.
\end{abstract}
\input{sections/introduction}
\input{sections/methods}
\input{sections/experiments}
\input{sections/discussion}
\clearpage
{\small
\bibliographystyle{plainnat}
\bibliography{references}
}
\clearpage
\appendix
\input{sections/appendix}
\clearpage
\input{sections/checklist}
\end{document}
